\usepackage[margin=2.35cm,headheight=15pt]{geometry}
\usepackage{amsmath,amssymb,booktabs,tabularx,longtable,array}
\usepackage{graphicx,xcolor,listings,enumitem,caption,placeins,fancyhdr,float}
\usepackage{hyperref}
\hypersetup{colorlinks=true,linkcolor=black,citecolor=black,urlcolor=blue,
  pdftitle={不同优先队列下 Dijkstra 算法的实现与评测},pdfauthor={},
  pdfsubject={项目二报告底稿与可复现实验},pdfcreator={XeLaTeX}}
\graphicspath{{figures/}}
\setlength{\parindent}{2em}
\setlength{\parskip}{0.2em}
\setlength{\emergencystretch}{3em}
\renewcommand{\arraystretch}{1.15}
\renewcommand{\tabularxcolumn}[1]{>{\raggedright\arraybackslash}p{#1}}
\raggedbottom
\setlist{itemsep=0.15em,topsep=0.3em}
\captionsetup{font=small,labelfont=bf}
\ctexset{section={name={第,章},number=\arabic{section}},subsection={format=\large\bfseries}}
\setcounter{tocdepth}{2}
\pagestyle{fancy}
\fancyhf{}
\fancyhead[L]{\small Dijkstra 算法与优先队列}
\fancyhead[R]{\small 项目二报告底稿}
\fancyfoot[C]{\thepage}
\renewcommand{\headrulewidth}{0.3pt}
\newcommand{\code}[1]{\texttt{\detokenize{#1}}}
\newcommand{\drafttodo}[1]{\ifDraftNotes\par\noindent\textbf{待补充：}#1\par\fi}
\newcommand{\plotfigure}[3]{%
  \begin{figure}[H]\centering
  \includegraphics[width=\linewidth]{#1.pdf}
  \caption{#2}\label{#3}
  \end{figure}}
\lstset{basicstyle=\ttfamily\small\linespread{1.0}\selectfont,columns=fullflexible,breaklines=true,
  keepspaces=true,showstringspaces=false,numbers=left,numberstyle=\tiny,
  frame=lines,tabsize=2,captionpos=b}
